\chapter{原子的壳层结构}
\problem{7.1}{有两种原子，在基态时其电子壳层是这样填充的：
\begin{enumerate}
\item $n=1$ 壳层、$n=2$ 壳层和 $3\mathrm s$ 次壳层都填满，$3\mathrm p$ 次壳层填了一半。
\item $n=1$ 壳层、$n=2$ 壳层、$n=3$ 壳层及 $4\mathrm s,4\mathrm p,4\mathrm d$ 次壳层都填满。
\end{enumerate}
试问这是哪两种原子？}
\answerarea
\problem{7.2}{原子的 $3\mathrm d$ 次壳层按泡利原理一共可以填多少电子？为什么？}
\answerarea\clearpage
\problem{7.3}{$\mathrm{Na}$ 原子的 $\mathrm S,\mathrm P,\mathrm D$ 项的量子修正值 $\Delta_{\mathrm S}=1.35,\Delta_{\mathrm P}=0.86,\Delta_{\mathrm D}=0.01$。把谱项表达成 $R(Z-\sigma)^2/n^2$ 形式，其中 $Z$ 是核电荷数。试计算 $3\mathrm S,3\mathrm P,3\mathrm D$ 项的 $\sigma$ 分别为何值？并说明 $\sigma$ 的物理意义。}
\answerarea
\problem{7.4}{原子中能够有下列量子数相同的最大电子数是多少？
\begin{enumerate}
\item $n,l,m_l$；
\item $n,l$；
\item $n$。
\end{enumerate}}
\answerarea\clearpage
\problem{7.5}{从实验得到的等电子体系 $\mathrm{K\,I},\mathrm{Ca\,II},\ldots$ 等的莫塞莱图解，怎样知道从钾 $Z=19$ 开始不填 $3\mathrm d$ 而填 $4\mathrm s$ 次壳层，又从钪 $Z=21$ 开始填 $3\mathrm d$ 而不填 $4\mathrm s$ 次壳层？}
\answerarea\clearpage
